%!TEX root = ../main.tex
%!TEX program = xelatex
% Chapter 7: Words 1501 - 1750
\chapter{Professional and Literary Vocabulary}
\label{chap:chapter7}

{\small\sffamily\color{mutedtext} Higher-register verbs, formal dialogue, culture, science, and societal terminology.\par}

\vspace{0.4em}
\markboth{Chapter 7: Professional and Literary Vocabulary}{Words 1501--1750}

\begin{multicols}{2}
\vocentry{1501}{Старшина}{starshi'na}{A sergeant major (a military rank)}{Старшина поприветствовал новобранцев.}{The sergeant major greeted the recruits.}
\vocentry{1502}{Озеро}{'ozerə}{A lake}{Горное озеро было чистым и прозрачным.}{The mountain lake was pure and transparent.}
\vocentry{1503}{Чемодан}{tchema'dan}{A suitcase}{Ты упаковала свой чемодан? Не забудь про тёплые вещи!}{Have you packed your suitcase? Don’t forget about the warm clothes!}
\vocentry{1504}{Вынуть}{'vynut’}{To take out, to pull out (colloquial)}{Что за странный запах? Ты забыл вынуть пирог из духовки?}{What kind of a strange smell is it? Have you forgotten to take the pie out of the oven?}
\vocentry{1505}{Сзади}{'zzadi}{Behind, from behind, in the back}{Пассажиры сзади меня громко разговаривали всю дорогу!}{The passengers behind me were talking loudly the whole way!}
\vocentry{1506}{Стоить}{'stoit’}{To cost; to be worth}{Эти туфли не могут столько стоить! Я не собираюсь их покупать!}{These shoes can’t cost so much! I’m not going to buy them!}
\vocentry{1507}{Предел}{pre'del}{A limit, a boundary, a margin}{У моего терпения есть предел. Я устала работать сверхурочно.}{There’s a limit to my patience. I’m tired of working overtime.}
\vocentry{1508}{Музей}{mu'zej}{A museum}{На выходных мы посетили морской музей.}{We visited a maritime museum at the weekend.}
\vocentry{1509}{Пятно}{pjat'no}{A stain, a spot}{Я не могу вывести это жирное пятно со своей любимой блузки.}{I can’t get this grease stain out of my favorite blouse.}
\vocentry{1510}{Зона}{'zona}{A zone; a prison (slang)}{Это радиоактивная зона. Здесь никто не живёт уже много лет.}{It’s a radiation zone. Nobody has been living here for many years.}
\vocentry{1511}{Ученик}{utche'nik}{A pupil, a student}{Работать с этим мальчиком одно удовольствие! Он очень способный ученик.}{Working with this boy is a sheer pleasure! He’s a very gifted pupil.}
\vocentry{1512}{Навсегда}{navseg'da}{Forever}{Я собирался пожить в Лондоне пару лет, а остался здесь навсегда.}{I was going to live in London for a couple of years but stayed here forever.}
\vocentry{1513}{Толк}{tolk}{A use, a sense}{Какой толк от твоих извинений? Они ничего не изменят.}{What’s the use of your apologies? They won’t change anything.}
\vocentry{1514}{Полагать}{pala'gat’}{To suppose, to believe, to recon}{У меня есть причина полагать, что подозреваемый не виноват.}{I’ve got a reason to suppose that the suspect is not guilty.}
\vocentry{1515}{Родственник}{'rotstvenik}{A relative}{Он мой дальний родственник, и я практически ничего не знаю о нём.}{He’s my distant relative and I know almost nothing about him.}
\vocentry{1516}{Источник}{is'tochnik}{A source; a spring}{Этот веб сайт – ненадёжный источник информации.}{This website is an unreliable source of information.}
\vocentry{1517}{Угодно}{u'godnə}{Most often used in a word combination ‘какой угодно’ - {\ipafont [ka'koj u'godn{\ipafont ə}]} - Any, whichever, whatsoever}{Давай, задай мне какой угодно вопрос! Я отвечу на всё!}{Come on, ask me any question! I’ll answer all of them!}
\vocentry{1518}{Послушать}{pas'lushat’}{To listen (to listen for some time)}{Я хочу немного отдохнуть, послушать музыку и расслабиться.}{I want to have some rest, to listen to some music and relax.}
\vocentry{1519}{Колесо}{kale'so}{A wheel}{Я не знаю, кто изобрёл колесо, но этот человек был гений.}{I don’t know who invented the wheel but that man was a genius.}
\vocentry{1520}{Грязь}{gr’as’}{Mud, filth, dirt}{На твоих ботинках грязь. Не забудь почистить их.}{There’s mud on your shoes. Don’t forget to clean them.}
\vocentry{1521}{Отметить}{at'metit’}{To celebrate; to point out; to mark}{Я хочу отметить Новый год в семейном кругу.}{I want to celebrate the New Year within the family circle.}
\vocentry{1522}{Стен}{sten}{A wall (in genitive plural of ‘стена’)}{Из-за сырости обои отклеились от стен.}{The wallpaper peeled off the walls because of dampness.}
\vocentry{1523}{Добро}{dab'ro}{The good}{Он всегда стремился делать людям добро, несмотря ни на что.}{He’s always strived to do people good, no matter what.}
\vocentry{1524}{Капля}{'kaplja}{A drop; a bit}{Говорят, что капля никотина может убить лошадь.}{They say a drop of nicotine can kill a horse.}
\vocentry{1525}{Экран}{ɛk'ran}{A screen}{Мне нужен телевизор побольше. У этого экран маленький.}{I need a bigger TV set. This one has a small screen.}
\vocentry{1526}{Руководитель}{rukava'ditel’}{A leader, a manager}{Настоящий лидер работает вместе со своей командой.}{A real leader works together with his team.}
\vocentry{1527}{Инспектор}{ins'pektər}{An inspector, a controller}{Инспектор остановил подозрительный автомобиль, который ехал слишком быстро.}{The inspector stopped a suspicious car that was going too fast.}
\vocentry{1528}{Удивление}{udiv'lenije}{Amazement, surprise}{Она не могла скрыть своё удивление, когда увидела, что мы убрали всю квартиру.}{She couldn’t hide her amazement when she saw that we’d cleaned the whole apartment.}
\vocentry{1529}{Утверждать}{utver'zhdat’}{To claim, to argue}{Как ты можешь утверждать, что она виновата, если ты даже не знаешь, что случилось?}{How can you claim that she’s guilty if you don’t even know what had happened?}
\vocentry{1530}{Молоко}{mala'ko}{Milk}{Горячее молоко с мёдом поможет облегчить твой кашель.}{Some hot milk with honey will help to relieve your cough.}
\vocentry{1531}{Очевидно}{atche'vidnə}{Obviously, it is obvious}{Вы, очевидно, не совсем понимаете, что мы от вас хотим.}{You obviously don’t quite understand what we want from you.}
\vocentry{1532}{Отличаться}{atli'tchatsa}{To differ}{Эти элементы не должны отличаться друг от друга по цвету и размеру.}{These elements shouldn’t differ from each other in color and size.}
\vocentry{1533}{Предложение}{predla'zhenije}{An offer; a sentence}{Это очень интересное деловое предложение, но мне придётся отказаться.}{It’s a very interesting business offer but I’ll have to decline.}
\vocentry{1534}{Нету}{'netu}{No, there is no, there are no (colloquial)}{У меня совсем нету времени на пустые разговоры.}{I have absolutely no time for empty talk.}
\vocentry{1535}{Среда}{sre'da}{Environment, surroundings; Wednesday}{Этим растениям нужна особая среда.}{These plants need a special environment.}
\vocentry{1536}{Выскочить}{'vyskatchit’}{To jump out, to pop out}{Будь осторожен! Здесь на дорогу могут выскочить дикие животные.}{Be careful! Wild animals can jump out onto the road here.}
\vocentry{1537}{Одетый}{a'detyj}{Dressed}{На вечеринке был парень, одетый как пришелец.}{At the party there was a guy dressed like an alien.}
\vocentry{1538}{Площадка}{pla'sthadka}{A platform; a playground}{На верху здания есть смотровая площадка.}{There’s a viewing platform on the top of the building.}
\vocentry{1539}{Приезжать}{priez'zhat’}{To come, to arrive (by horse or vehicle)}{Я люблю приезжать в родительский дом во время каникул.}{I like to come to my parents’ house during the holidays.}
\vocentry{1540}{Катить}{ka'tit’}{To roll, to push (an object with wheels)}{В этом чемодане есть колёсики. Ты можешь катить его.}{This suitcase has wheels. You can roll it.}
\vocentry{1541}{Совесть}{'sovest’}{Conscience}{Если у тебя есть сомнения, слушай свою совесть.}{If you have doubts, listen to your conscience.}
\vocentry{1542}{Кофе}{'kofe}{Coffee}{Я не люблю кофе с сахаром. Он только портит вкус.}{I don’t like coffee with sugar. It only spoils the taste.}
\vocentry{1543}{Штаны}{shta'ny’}{Pants, trousers (more for casual style)}{Эти штаны продаются со скидкой. Хочешь примерить их?}{These pants are on discount. Want to try them on?}
\vocentry{1544}{Охота}{a'hota}{Hunting}{Я считаю, что охота – это жестокое и бессмысленное развлечение.}{I consider hunting to be cruel and senseless entertainment.}
\vocentry{1545}{Бровь}{brov’}{An eyebrow}{Оса ужалила меня прямо в левую бровь.}{A wasp stung me right on the left eyebrow.}
\vocentry{1546}{Расти}{ras'ti}{To grow; to increase}{Эти деревья могут расти только в тропическом климате.}{These trees can only grow in a tropical climate.}
\vocentry{1547}{Список}{'spisək}{A list}{Проверь список гостей. Может, ты захочешь добавить кого-нибудь.}{Check the guest list. Maybe you’ll want to add someone.}
\vocentry{1548}{Ниже}{'nizhə}{Below; lower}{Ниже вы найдёте все прикреплённые документы.}{Below you’ll find all the attached documents.}
\vocentry{1549}{Вернуть}{ver'nut’}{To return, to give back}{Я постараюсь вернуть твою книгу как можно скорее.}{I’ll try to return your book as soon as possible.}
\vocentry{1550}{Зависеть}{za'viset’}{To depend on}{Когда я была студенткой, я не хотела зависеть от родителей и нашла себе работу.}{When I was a student, I didn’t want to depend on my parents and found myself a job.}
\vocentry{1551}{Ствол}{stvol}{A trunk (of a tree); a gun (slang)}{Ствол этого дерева просто огромный! Я не могу обхватить его руками.}{The trunk of this tree is just huge! I can’t put my arms around it.}
\vocentry{1552}{Основание}{asna'vanije}{A ground, a reason; a basis}{У тебя есть хоть какое-нибудь основание не доверять мне?}{Do you have at least any ground not to trust me?}
\vocentry{1553}{Северный}{'severnyj}{Northern}{Северный район города сильно пострадал от вчерашнего града.}{The northern district of the city suffered from yesterday’s hail a lot.}
\vocentry{1554}{Познакомиться}{pazna'komitsa}{To meet, to get acquainted}{Приятно познакомиться с вами! Я много слышал о вас от своей сестры.}{Nice to meet you! I’ve heard a lot about you from my sister.}
\vocentry{1555}{Поведение}{pave'denije}{Behavior}{Такое поведение просто неприемлемо! Нельзя вот так выходить из себя.}{Such behavior is just unacceptable! One can’t lose their temper like that.}
\vocentry{1556}{Жалить}{'zhalit’}{To sting (of insects); to bite (of snakes)}{Пчела не будет жалить, если ты не будешь её трогать.}{A bee won’t sting you if you don’t touch it.}
\vocentry{1557}{Найтись}{naj'tis’}{To be found, to be discovered (in passive)}{Смотри, вот твоя серёжка! Она должна была найтись рано или поздно.}{Look, here’s your earring! It had to be found sooner or later.}
\vocentry{1558}{Делаться}{'delatsa}{To be done, to be made}{Такая работа не должна делаться в спешке.}{Such work mustn’t be done in haste.}
\vocentry{1559}{Догадаться}{daga'datsa}{To guess}{Он не должен догадаться, что мы готовим ему сюрприз.}{He mustn’t guess that we’re planning a surprise for him.}
\vocentry{1560}{Дальнейший}{dal’'nejshij}{Further}{Дальнейший прогресс твоего выздоровления зависит только от тебя.}{Further progress of your recovery depends only on you.}
\vocentry{1561}{Прожить}{pra'zhit’}{To live (for some time or to live a life)}{Я хочу прожить долгую жизнь, поэтому забочусь о своём здоровье.}{I want to live a long life which is why I care about my health.}
\vocentry{1562}{Усилие}{u'silije}{An effort}{Ещё одно усилие – и мы сдвинем этот камень.}{One more effort, and we’ll move this stone.}
\vocentry{1563}{Гулять}{gu'ljat’}{To walk, to go for a walk, to stroll}{Я люблю гулять по ночному городу.}{I like to walk around the city at night.}
\vocentry{1564}{Клетка}{'kletka}{A cage; a cell (biology)}{Мне нужна более просторная клетка для моего попугая.}{I need a more spacious cage for my parrot.}
\vocentry{1565}{Разобраться}{razab'ratsa}{To sort out; to understand}{Я не смогу разобраться с этой проблемой в одиночку. Мне нужна помощь.}{I won’t be able to sort this problem out alone. I need help.}
\vocentry{1566}{Обращать}{abrast'chat’}{To pay (most often about attention)}{Постарайся не обращать внимания на его поведение. Он делает это нарочно.}{Try not to pay attention to his behavior. He does it on purpose.}
\vocentry{1567}{Верхний}{'verhnij}{Upper}{Верхний слой воды в озере тёплый, но глубже вода ледяная.}{The upper layer of the water in the lake is warm but deeper down the water is icy cold.}
\vocentry{1568}{Обстановка}{absta'novka}{An atmosphere, a setting; decor, furniture}{После ссоры обстановка в комнате была напряжённой.}{After the quarrel the atmosphere in the room was tense.}
\vocentry{1569}{Температура}{tempera'tura}{Temperature}{Температура кипящей воды – 100 градусов по Цельсию.}{The temperature of boiling water is 100 degrees Celcius.}
\vocentry{1570}{Вывод}{'vyvət}{A conclusion; withdrawal}{Не спеши делать вывод. Ты ещё не знаешь всех обстоятельств.}{Don’t hurry to make a conclusion. You don’t know all the circumstances yet.}
\vocentry{1571}{Абсолютно}{absa'ljutnə}{Absolutely, completely}{Я абсолютно уверен, что они помирятся.}{I’m absolutely sure that they’ll make up.}
\vocentry{1572}{Предприятие}{pretpri'jatije}{An enterprise, a business}{Это семейное предприятие. Оно было открыто моим дедом семьдесят лет назад.}{It’s a family enterprise. It was opened by my grandfather seventy years ago.}
\vocentry{1573}{Принятый}{'prinjatyj}{Adopted, accepted}{Мне кажется, что принятый вчера закон поможет развитию малого бизнеса.}{It seems to me that the law adopted yesterday will help the development of the small business.}
\vocentry{1574}{Отдавать}{atda'vat’}{To give back, to give, to return}{Девочка не хотела отдавать игрушку, которую взяла у друга в песочнице.}{The girl didn’t want to give back the toy she borrowed from a friend in the sandbox.}
\vocentry{1575}{Сходить}{sha'dit’}{To go (perfective) (means to go and do something)}{Ты можешь сходить на кухню и принести печенье?}{Can you go to the kitchen and bring the cookies?}
\vocentry{1576}{Подарок}{pa'darək}{A gift, a present}{Сложно найти подарок для человека, у которого всё есть!}{It’s hard to find a gift for a person who has everything!}
\vocentry{1577}{Звучать}{zvu'tchat’}{To sound}{Твои слова будут звучать убедительнее, если ты перестанешь кричать.}{Your words will sound more persuasive if you stop yelling.}
\vocentry{1578}{Отделение}{atde'lenije}{A unit, a department; a compartment}{Его отвезли в отделение реанимации. Надеюсь, врачам удастся спасти его жизнь.}{They took him to the intensive care unit. Hopefully, the doctors will manage to save his life.}
\vocentry{1579}{Волноваться}{valna'vatsa}{To worry}{Не нужно так волноваться. Есть много причин, по которым она может не брать трубку.}{There’s no need to worry so much. There are lots of reasons why she may not pick up the phone.}
\vocentry{1580}{Нижний}{'nizhnij}{Lower, inferior}{На этой фотографии есть наш дедушка. Смотри, нижний ряд, первый справа.}{There’s our grandfather in this photo. Look, the lower row, the first on the right.}
\vocentry{1581}{Штука}{'shtuka}{A thing (colloquial); a piece, an item}{Что это за странная штука? И почему она лежит у меня на кровати?}{What kind of a strange thing is it? And why is it lying on my bed?}
\vocentry{1582}{Защита}{za'sthita}{Protection, defense}{Защита вымирающих видов – это дело всей его жизни.}{The protection of endangered species is a life work of his.}
\vocentry{1583}{Вряд ли}{'vrjadli}{Unlikely, not likely, hardly}{Вряд ли это что-то серьёзное. Я думаю, ты просто простудился.}{It’s unlikely it is something serious. I think you’ve just caught a cold.}
\vocentry{1584}{Выстрел}{'vystrel}{A shot, a gunshot}{Пожилая женщина сказала следователю, что слышала выстрел около семи часов вечера.}{The elderly lady told the detective that she’d heard the shot at about seven p.m.}
\vocentry{1585}{Тотчас}{'totchas}{Immediately, at once (literary style)}{Мы тотчас отправляем помощь! Оставайтесь на месте!}{We’re sending help immediately! Stay where you are!}
\vocentry{1586}{Суть}{sut’}{Essence, gist}{Суть твоей идеи мне ясна. Но как мы реализуем её?}{The essence of your idea is clear to me. But how are we going to realize it?}
\vocentry{1587}{Труп}{trup}{A corpse, a dead body}{Детективы нашли труп в лесу, недалеко от дороги.}{The detectives found the corpse in the forest, not far from the road.}
\vocentry{1588}{Основа}{as'nova}{A basis, a foundation}{Любовь и понимание – это основа любых отношений.}{Love and understanding is the basis of any relationship.}
\vocentry{1589}{Подать}{pa'dat’}{To give (perfective) (means to give something to someone who’s far from the object)}{Можешь подать мне ту большую книгу, которая рядом с тобой?}{Can you give me the big book that’s next to you?}
\vocentry{1590}{Домашний}{da'mashnij}{Home-made; home (adj)}{Я удивлён! Этот консервированный суп на вкус почти как домашний.}{I’m surprised! This canned soup tastes almost like a home-made one.}
\vocentry{1591}{Готовить}{ga'tovit’}{To cook; to prepare}{Обычно я люблю готовить, но сегодня мне лень и я заказала пиццу.}{Usually I like to cook but I’m lazy today and I’ve ordered some pizza.}
\vocentry{1592}{Участие}{u'tchastije}{Participation}{Я не одобряю твоё участие в этом музыкальном конкурсе. Какая в нём польза?}{I disapprove of your participation in this music contest. What’s the use of it?}
\vocentry{1593}{Жалеть}{zha'let’}{To regret, to be sorry}{Если ты не попробуешь сейчас, ты будешь жалеть об этом всю свою жизнь.}{If you don’t try now, you’ll regret it all your life.}
\vocentry{1594}{Фон}{fon}{A background}{Для этого яркого изображения мне нужен более нейтральный фон.}{I need a more neutral background for this bright image.}
\vocentry{1595}{Признак}{'priznak}{A sign, an indication}{Доктор сказал, что хороший аппетит – это признак выздоровления.}{The doctor said that good appetite is a sign of recovery.}
\vocentry{1596}{Врать}{vrat’}{To lie}{Не буду врать, мне всегда хотелось быть такой, как ты.}{I won’t lie, I’ve always wanted to be like you.}
\vocentry{1597}{Возразить}{vazra'zit’}{To object, to mind, to protest}{Девушка хотела возразить, но шум заглушил её голос.}{The girl wanted to object but the noise silenced her voice.}
\vocentry{1598}{Дача}{'datcha}{A country house}{У нас есть дача, но она далеко от города.}{We have a country house but it’s far from the city.}
\vocentry{1599}{Обо}{'obə}{About (a variation of preposition ‘о’ used before consonant clusters)}{Ты ничего не знаешь обо мне!}{You know nothing about me!}
\vocentry{1600}{Довольный}{da'vol’nyj}{Satisfied, pleased}{После ужина довольный кот мирно уснул у меня на коленях.}{After supper the satisfied cat peacefully fell asleep on my lap.}
\vocentry{1601}{Готовиться}{ga'tovitsa}{To get ready, to prepare oneself}{Завтра я буду готовиться к экзаменам и не смогу пойти с тобой по магазинам.}{I’m going to get ready for my exams tomorrow and won’t be able to go shopping with you.}
\vocentry{1602}{Слой}{sloj}{A layer}{Перед покраской мы удалили толстый слой ржавчины.}{Before painting we removed a thick layer of rust.}
\vocentry{1603}{Охрана}{ah'rana}{Security, guard, protection}{Охрана супермаркета задержала магазинного воришку.}{The supermarket security detained a shoplifter.}
\vocentry{1604}{Превратиться}{prevra'titsa}{To turn into (intransitive)}{Иногда я хочу превратиться в птицу и улететь далеко отсюда.}{Sometimes I want to turn into a bird and fly far away from here.}
\vocentry{1605}{Милиционер}{militsia'ner}{A policeman, a militiaman}{За спасение ребёнка милиционер был награжден медалью.}{The policeman was awarded a medal for saving a child.}
\vocentry{1606}{Портрет}{par'tret}{A portrait}{Этот прекрасный портрет выставлен в музее, но никто не знает, кто его написал.}{This wonderful portrait is displayed in the museum but nobody knows who painted it.}
\vocentry{1607}{Терпеть}{ter'pet’}{To endure, to stand, to tolerate}{Зачем терпеть боль, если ты можешь принять таблетку?}{Why endure the pain if you can take a pill?}
\vocentry{1608}{Махнуть}{mah'nut’}{To wave; to swap (colloquial)}{Я попросил его махнуть рукой, когда он будет готов.}{I asked him to wave his hand when he’s ready.}
\vocentry{1609}{Шкаф}{shkaf}{A wardrobe, a closet}{Если ты продолжишь покупать столько одежды, нам придётся купить ещё один шкаф.}{If you go on buying so many clothes, we’ll have to buy one more wardrobe.}
\vocentry{1610}{Вес}{ves}{Weight}{Избыточный вес опасен для сердца.}{Excessive weight is dangerous for the heart.}
\vocentry{1611}{Холод}{'holəd}{Cold (low temperature)}{Я люблю зимний холод. Если он умеренный, конечно.}{I like winter cold. If it’s moderate, of course.}
\vocentry{1612}{Определить}{aprede'lit’}{To determine, to define, to identify (perfective)}{Сложно определить точное количество жертв землетрясения.}{It’s hard to determine the exact number of earthquake victims.}
\vocentry{1613}{Выпустить}{'vypustit’}{To release; to produce, to issue}{Постарайся выпустить плохую энергию и сконцентрируйся на хорошем.}{Try to release bad energy and concentrate on the good things.}
\vocentry{1614}{Тревога}{tre'voga}{Alarm, alert; anxiety}{Прошлой ночью пожарная тревога разбудила весь район, но она спасла много жизней.}{Last night the fire alarm woke the whole neighborhood up but it saved lots of lives.}
\vocentry{1615}{Выступать}{vystu'pat’}{To perform (on the stage)}{Завтра моя любимая группа будет выступать в центральном концертном зале.}{Tomorrow my favorite band will perform in the central concert hall.}
\vocentry{1616}{Полок}{'polək}{Shelves (genitive plural)}{Зачем тебе столько книжных полок? Насколько я знаю, ты не большой фанат чтения.}{Why do you need so many book shelves? As far as I know, you’re not a great fan of reading.}
\vocentry{1617}{Куртка}{'kurtka}{A jacket}{Тебе нужна куртка потеплее. Эта слишком тонкая.}{You need a warmer jacket. This one is too thin.}
\vocentry{1618}{Спустя}{spus'tja}{Later, after}{Два года спустя я вернулся на то место, где посадил молодые деревья.}{Two years later I returned to the place where I’d planted young trees.}
\vocentry{1619}{Отчего}{atche'vo}{Why (because of what reason; more often in literary style)}{Отчего ты такой грустный сегодня? Что-то не так на работе?}{Why are you so sad today? Anything wrong at work?}
\vocentry{1620}{Изба}{iz'ba}{A peasant’s hut (archaic, used in modern language to ironically call a house or in set expressions like the one in the example that is literally translated as ‘To take out dirt from one’s house)}{Зачем выносить сор из избы? Давай обсудим всё лично.}{Why wash our dirty linen in public? Let’s discuss everything in private.}
\vocentry{1621}{Редко}{'retkə}{Rarely, seldom}{Я редко ем фастфуд, но вчера решила себя побаловать.}{I rarely eat fast food but yesterday I decided to indulge myself.}
\vocentry{1622}{Вершина}{ver'shina}{A peak, a summit}{Эта горная вершина не тает зимой, и она очень популярна среди туристов.}{This mountain peak doesn’t melt in winter and is very popular among tourists.}
\vocentry{1623}{Тайный}{'tajnyj}{Secret (adj)}{Историки говорят, что под зданием есть тайный ход.}{Historians say that there’s a secret passage under the building.}
\vocentry{1624}{Занятый}{'zanjatyj}{Engaged; occupied}{Водитель, занятый телефонным разговором, не заметил пешехода.}{The driver engaged in a phone call didn’t notice the pedestrian.}
\vocentry{1625}{Рассматривать}{rass'matrivat’}{To consider, to address}{Неправильно рассматривать конфликт только с одной точки зрения.}{It’s not right to consider a conflict from only one point of view.}
\vocentry{1626}{Отойти}{ataj'ti}{To move aside (perfective)}{Я выхожу на следующей остановке. Могу я попросить вас отойти немного?}{I’m getting off at the next stop. May I ask you to move aside a bit?}
\vocentry{1627}{Вырасти}{'vyrasti}{To grow, to increase (perfective)}{Эта должность дала мне возможность вырасти профессионально.}{That position gave me an opportunity to grow professionally.}
\vocentry{1628}{Париж}{pa'rizh}{Paris}{На наш медовый месяц мы поехали в Париж.}{We went to Paris on our honeymoon.}
\vocentry{1629}{Убийство}{u'bijstvə}{A murder}{Никто не знает, кто совершил это убийство, но есть несколько подозреваемых.}{Nobody knows who committed this murder but there are a few suspects.}
\vocentry{1630}{Рубашка}{ru'bashka}{A shirt}{Зачем тебе ещё одна рубашка? Мне кажется, у тебя их тысячи.}{Why do you need another shirt? It seems to me you’ve got thousands of them.}
\vocentry{1631}{Ох}{oh}{Oh (mostly to express sadness or tiredness)}{Ох, как я устала! А сегодня ещё только понедельник.}{Oh, how tired I am! And it’s only Monday today.}
\vocentry{1632}{Хозяйство}{ha'zjastvə}{A household; economy; a farm}{Вести хозяйство и воспитывать детей не так просто, как кажется.}{Maintaining the household and bringing up children is not as easy as it seems.}
\vocentry{1633}{Парк}{park}{A park}{Парк вокруг дворца знаменит редкими видами деревьев и цветов.}{The park around the palace is famous for the rare species of trees and flowers.}
\vocentry{1634}{Перевести}{pereves'ti}{To translate, to interpret; to transfer (perfective)}{Я не смогу перевести этот текст без словаря.}{I won’t be able to translate this text without a dictionary.}
\vocentry{1635}{Режим}{re'zhim}{A mode; a regime}{Переключись на спящий режим, если хочешь сэкономить заряд батареи.}{Turn to sleep mode if you want to save the battery power.}
\vocentry{1636}{Руководство}{ruka'vodstvə}{A guide, an instruction; leadership}{Перед использованием прибора внимательно прочитайте руководство пользователя.}{Before using the appliance read the user guide attentively.}
\vocentry{1637}{Посадить}{pasa'dit’}{To plant; to put to jail}{Я хочу посадить перед домом много деревьев, чтобы у меня был свой сад.}{I want to plant lots of trees in front of my house to have my own garden.}
\vocentry{1638}{Преступление}{prestup'lenije}{A crime}{Это преступление заслуживает самого жестокого наказания.}{This crime deserves the most severe punishment.}
\vocentry{1639}{Стенка}{'stenka}{A side (of a container); a wall (diminutive of ‘стена’);}{Нижняя стенка резервуара должна быть сделана из водонепроницаемого материала.}{The lower side of the reservoir must be made of a waterproof material.}
\vocentry{1640}{Орать}{a'rat’}{To yell, to scream}{Зачем так орать? Успокойся и объясни всё ещё раз.}{Why yell so much? Calm down and explain everything once again.}
\vocentry{1641}{Кругом}{kru'gom}{Around, round}{Кругом было так много сладостей и угощений, что я не знала, что выбрать.}{There were so many sweets and treats around that I didn't know what to choose.}
\vocentry{1642}{Обойтись}{abaj'tis’}{To do without (perfective)}{Мы не сможем обойтись без дополнительной помощи. У нас слишком мало опыта.}{We won't be able to do it without extra help. We've got too little experience.}
\vocentry{1643}{Луна}{lu'na}{The moon}{Луна спряталась за облаками, и стало совсем темно.}{The moon hid behind the clouds and it got completely dark.}
\vocentry{1644}{Богатый}{ba'gatyj}{Rich; a rich man}{Неужели он настолько богатый, что у него есть собственный самолёт?}{Is he really so rich that he has his own airplane?}
\vocentry{1645}{Гостиница}{gas'tinitsa}{A hotel, an inn}{Мне нравится эта гостиница: хороший сервис, доступная цена.}{I like this hotel: good service, an affordable price.}
\vocentry{1646}{Беседа}{be'seda}{A talk, a conversation}{Иногда всё, что нужно, – это душевная беседа за чашечкой чая.}{Sometimes all you need is a soulful talk over a cup of tea.}
\vocentry{1647}{Твёрдый}{'tvjordyj}{Firm, solid}{В темноте я споткнулась о какой-то твёрдый предмет и упала.}{In the darkness I stumbled over some firm object and fell.}
\vocentry{1648}{Полк}{polk}{A regiment (military)}{Когда солдат восстановился после ранения, он вернулся в полк.}{When the soldier recovered from the wound he returned to the regiment.}
\vocentry{1649}{Выдать}{'vydat’}{To give out}{Как ты мог выдать мой секрет? Ты обещал, что никому не скажешь!}{How could you give out my secret? You promised you wouldn’t tell anyone!}
\vocentry{1650}{Лапа}{'lapa}{A paw}{Лапа этого огромного пса больше твоей руки!}{The paw of this huge dog is bigger than your hand!}
\vocentry{1651}{Воевать}{vaje'vat’}{To be at war, to wage war}{Я не хочу воевать с тобой. Давай попробуем помириться.}{I don’t want to be at war with you. Let’s try to make up.}
\vocentry{1652}{Останавливаться}{asta'navlivatsa}{To stop (reflexive)}{Мы уже сделали так много, и глупо останавливаться сейчас.}{We’ve already done so much, and it’s silly to stop now.}
\vocentry{1653}{Никуда}{niku'da}{(To) nowhere}{У неё ужасная депрессия. Она никуда не ходит и ни с кем не общается.}{She’s got a terrible depression. She goes nowhere and doesn’t communicate with anyone.}
\vocentry{1654}{Обладать}{abla'dat’}{To possess, to own}{Нужно обладать огромным терпением, чтобы жить с таким человеком, как твой муж.}{One must possess enormous patience to live with a person like your husband.}
\vocentry{1655}{Продать}{pra'dat’}{To sell (perfective)}{Извините, но я не могу продать эти часы. Это семейная реликвия.}{Sorry, but I can’t sell this watch. It’s a family heirloom.}
\vocentry{1656}{Спорить}{'sporit’}{To dispute, to argue}{Даже не пытайся спорить с ним. Его невозможно переубедить.}{Don’t even try to dispute with him. It’s impossible to dissuade him.}
\vocentry{1657}{Солнечный}{'solnetchnyj}{Sunny}{Жаль, что тебе приходится оставаться дома в такой солнечный день.}{It’s a pity you have to stay at home on such a sunny day.}
\vocentry{1658}{Забор}{za'bor}{A fence}{Между нашими домами стоит высокий забор. Жизнь соседей меня совсем не интересует.}{There’s a high fence between our houses. I’m not interested in my neighbors’ life at all.}
\vocentry{1659}{Больший}{'bol’shij}{Greater, bigger, larger}{Мой сын проявляет больший интерес к книгам, чем к спорту.}{My son shows greater interest in books than in sports.}
\vocentry{1660}{Пост}{post}{A post, a position}{Маргарет Тэтчер занимала пост премьер-министра двенадцать лет.}{Margaret Thatcher held the post of Prime Minister for twelve years.}
\vocentry{1661}{Полезть}{pa'lest’}{To climb (perfective)}{Нужно быть сумасшедшим, чтобы полезть на такую высокую гору в одиночку.}{One must be crazy to climb such a high mountain alone.}
\vocentry{1662}{Печальный}{pe'tchal’nyj}{Sorrowful, grieved}{В последнее время у неё такой печальный вид, будто кто-то умер.}{She’s had such a sorrowful look recently, as if someone died.}
\vocentry{1663}{Выполнять}{vypal'njat’}{To carry out, to fulfill}{Пока моя коллега в отпуске, я буду выполнять её обязанности.}{While my colleague is on holiday I’ll be carrying out her duties.}
\vocentry{1664}{Требоваться}{'trebavatsa}{To be required, to be needed}{Для должной организации мероприятия будет требоваться много усилий.}{Much effort will be required for proper organisation of the event.}
\vocentry{1665}{Бледный}{'blednyj}{Pale}{Почему ты такой бледный? Ты что, увидел привидение?}{Why are you so pale? Have you seen a ghost?}
\vocentry{1666}{Умирать}{umi'rat’}{To die}{Никто не хочет умирать, особенно в таком молодом возрасте.}{Nobody wants to die, especially at such a young age.}
\vocentry{1667}{Шапка}{'shapka}{A hat (a warm piece of clothing, without brims)}{Эта шерстяная шапка тёплая, но очень колючая.}{This woolen hat is warm but very itchy.}
\vocentry{1668}{Дойти}{daj'ti}{To walk to, to reach}{До вокзала можно дойти за десять минут. Зачем садиться на автобус?}{One can walk to the railway station in ten minutes. Why take a bus?}
\vocentry{1669}{Успокоиться}{uspa'koitsa}{To calm down}{Мне нужно немного побыть одной, чтобы успокоиться. Оставьте меня.}{I need to be alone for some time to calm down. Leave me.}
\vocentry{1670}{Университет}{universi'tet}{A university}{В университет поступить не просто, но это того стоит.}{It’s not easy to enter a university but it’s worth it.}
\vocentry{1671}{Бабка}{'babka}{An old woman (colloquial, mostly scornfully)}{Та бабка, которая живёт на втором этаже, снова пыталась учить меня жизни!}{That old woman who lives on the second floor tried to teach me about life again!}
\vocentry{1672}{Журналист}{zhurna'list}{A journalist}{Этот журналист имеет своё мнение обо всех событиях. Я люблю читать его колонку.}{This journalist has his own opinion about all the events. I like reading his column.}
\vocentry{1673}{Должность}{'dolzhnast’}{A position, a post}{Он так сильно хотел эту должность, что был готов врать и унижаться.}{He wanted that position so much that he was ready to lie and grovel.}
\vocentry{1674}{Приближаться}{pribli'zhatsa}{To approach, to near}{Разбуди меня, когда мы начнём приближаться к нашей станции.}{Wake me up when we start approaching our station.}
\vocentry{1675}{Истина}{'istina}{Verity, truth}{Правда и истина – это разные вещи. Ты согласен?}{Truth and verity are different things. Do you agree?}
\vocentry{1676}{Раствор}{rast'vor}{Solution}{Я оставила свой раствор для контактных линз дома. Ты можешь одолжить мне свой?}{I left my contact lens solution at home. Can you lend me yours?}
\vocentry{1677}{Газ}{gas}{Gas}{Этот газ очень токсичен. Обращайтесь с ним осторожно.}{This gas is very toxic. Handle it with care.}
\vocentry{1678}{Позиция}{pa'zitsija}{A position, an attitude}{Твоя позиция по этому вопросу мне не ясна. Почему ты против?}{Your position on this matter is not clear to me. Why are you against it?}
\vocentry{1679}{Летать}{le'tat’}{To fly}{Этот птенец ещё не умеет летать и полностью зависит от родителей.}{This nestling can’t fly yet and depends entirely on its parents.}
\vocentry{1680}{Неожиданный}{nea'zhydannyj}{Unexpected, sudden}{Её неожиданный приезд всех нас приятно удивил.}{Her unexpected arrival pleasantly surprised us all.}
\vocentry{1681}{Убивать}{ubi'vat’}{To kill, to murder}{Курение продолжает убивать тысячи людей по всему миру.}{Smoking is continuing to kill thousands of people all over the world.}
\vocentry{1682}{Появление}{pajav'lenije}{Appearance, emergence}{Я никогда не забуду своё первое появление в новостях.}{I’ll never forget my first appearance in the news.}
\vocentry{1683}{Глухой}{glu'hoj}{Deaf}{Я не глухой. Не обязательно говорить так громко.}{I’m not deaf. It’s not necessary to talk so loudly.}
\vocentry{1684}{Обыкновенный}{abykna'vennyj}{Usual, ordinary}{Это был обыкновенный скучный рабочий день и я не ожидал никаких сюрпризов.}{That was an ordinary boring working day and I wasn’t expecting any surprises.}
\vocentry{1685}{Двинуться}{'dvinutsa}{To move (intransitive, perfective, means to start moving)}{Я так устал, что не могу двинуться с места.}{I’m so tired that I can’t move.}
\vocentry{1686}{Решительный}{re'shitel’nyj}{Decisive, resolute}{Что мне больше всего нравится в нём, так это то, что он решительный и никогда не сомневается.}{What I like about him most of all is that he’s decisive and never has doubts.}
\vocentry{1687}{Спуститься}{spust'titsa}{To go down, to decsend (perfective)}{Она обещала спуститься к ужину. Сходи проверь, что она делает.}{She promised to come down for supper. Go check what she’s doing.}
\vocentry{1688}{Песнь}{pesn’}{A song (archaic or poetic)}{О, песнь соловья так романтична!}{Oh, the nightingale's song is so romantic!}
\vocentry{1689}{Облако}{'oblakə}{A cloud}{Это белое облако напоминает котёнка, играющего с мячиком.}{This white cloud resembles a kitten playing with a ball.}
\vocentry{1690}{Морда}{'morda}{A muzzle; a face (colloquial, rude or ironic)}{Морда твоей собаки вся покрыта паутиной. Где она была?}{Your dog’s muzzle is all covered in spider web. Where has it been?}
\vocentry{1691}{Записка}{za'piska}{A note}{Где записка, которую оставила мама? Там написано, что мы должны сделать по дому.}{Where’s the note that mom left? It says what we must do around the house.}
\vocentry{1692}{Иностранный}{ina'strannyj}{Foreign}{Изучая иностранный язык, также важно узнавать об особенностях культуры.}{While studying a foreign language it’s also important to learn the peculiarities of the culture.}
\vocentry{1693}{Незнакомый}{nezna'komyj}{Unfamiliar, unknown}{Заходил какой-то незнакомый человек и сказал, что хочет поговорить с тобой.}{Some unfamiliar man dropped by and said he wanted to talk to you.}
\vocentry{1694}{Установить}{ustana'vit’}{To install, to mount}{Я скачал новую программу, но не знаю, как её установить.}{I’ve downloaded a new software program but I don’t know how to install it.}
\vocentry{1695}{Забывать}{zaby'vat’}{To forget}{Мы не должны забывать тех, кто отдал свои жизни ради нашей свободы.}{We shouldn’t forget those who gave their lives for the sake of our freedom.}
\vocentry{1696}{Позвать}{paz'vat’}{To call; to invite}{Можешь позвать сестру и сказать, что завтрак готов?}{Can you call your sister and tell her the breakfast is ready?}
\vocentry{1697}{Физический}{fi'zitcheskij}{Physical, bodily; physics (adj)}{Иногда психологический ущерб может иметь более серьёзные последствия, чем физический.}{Sometimes moral harm can have more serious consequences than physical harm.}
\vocentry{1698}{Ус}{us}{Most often used in plural ‘Усы - {\ipafont [u'sy]}’ - Moustache}{Сбрей свои усы. С ними ты выглядишь гораздо старше.}{Shave your moustache off. You look much older with it.}
\vocentry{1699}{Помещение}{pame'stchenije}{A room, a premise (as an opposite to outdoors)}{Не забывай снимать шляпу, когды входишь в помещение.}{Don’t forget to take your hat off when entering a room.}
\vocentry{1700}{Испугаться}{ispu'gatsa}{To get frightened, to get scared (perfective)}{Твой костюм Дракулы выглядит так реалистично, что можно испугаться.}{Your Dracula costume looks so realistic that one can get frightened.}
\vocentry{1701}{Плащ}{plastch}{A raincoat, an overcoat}{Этот плащ идеален для дождливой погоды. Он водонепроницаемый и достаточно тёплый.}{This raincoat is perfect for rainy weather. It’s waterproof and warm enough.}
\vocentry{1702}{Еле}{'jele}{Barely, hardly, scarcely}{Когда спасатели вытащили его из воды, он еле дышал.}{When the rescuers pulled him out of the water he was barely breathing.}
\vocentry{1703}{Темно}{tem'no}{Dark (adv)}{Уже становится темно. Нам лучше пойти домой, иначе мама будет волноваться.}{It’s getting dark already. We’d better go home otherwise Mom will be worried.}
\vocentry{1704}{Подъезд}{pad'jezd}{Entrance (the main one, most often in a block of flats)}{Мой подъезд первый от супермаркета. Номер квартиры ты помнишь?}{My entrance is the first one from the supermarket. Do you remember the apartment number?}
\vocentry{1705}{Съесть}{sjest’}{To eat, to eat up (perfective)}{Я такой голодный, что могу съесть целого слона!}{I’m so hungry I could eat a whole elephant!}
\vocentry{1706}{Прочесть}{pra'tchest’}{To read (completely), to read through (perfective)}{Я не смогла прочесть все книги из списка, который учитель задал нам на лето.}{I didn’t manage to read all the books from the list the teacher assigned us for summer.}
\vocentry{1707}{Следователь}{'sledavatel’}{An investigator, a detective}{Следователь допросил всех свидетелей и записал все детали в блокнот.}{The investigator interrogated all the witnesses and put down all the details in a notebook.}
\vocentry{1708}{Двенадцать}{dve'nadtsat’}{Twelve}{Когда часы пробьют двенадцать, карета снова превратится в тыкву.}{When the clock strikes twelve the carriage will turn into a pumpkin again.}
\vocentry{1709}{Задуматься}{za'dumatsa}{To think (means to begin to think), to reflect}{Тебе пора задуматься о том, чтобы завести детей.}{It’s time for you to think about having children.}
\vocentry{1710}{Решать}{re'shat’}{To solve, to sort out}{Ты не обязана решать эту проблему в одиночку. У тебя есть мы.}{You don’t have to solve this problem alone. You have us.}
\vocentry{1711}{Особенный}{a'sobennyj}{Special}{Я хочу найти для неё особенный подарок. Она заслуживает самого лучшего.}{I want to find a special present for her. She deserves the best.}
\vocentry{1712}{Подарить}{pada'rit’}{To give, to present}{Я не знаю, что подарить родителям на годовщину свадьбы.}{I don’t know what to give my parents for their wedding anniversary.}
\vocentry{1713}{Измениться}{izme'nitsa}{To change (intransitive, perfective)}{Мы планируем приехать в августе, но наши планы могут измениться.}{We’re planning to come in August, but our plans can change.}
\vocentry{1714}{Платок}{pla'tok}{A handkerchief; a shawl}{Я положила чистый платок в правый карман твоих брюк.}{I put a clean handkerchief into the right pocket of your trousers.}
\vocentry{1715}{Мощный}{'mostchnyj}{Powerful}{Мощный взрыв уничтожил практически всё здание.}{The powerful explosion destroyed almost all the building.}
\vocentry{1716}{Середина}{sere'dina}{A middle}{Сейчас середина месяца, а я уже потратил все деньги.}{It’s the middle of the month now and I’ve already spent all the money.}
\vocentry{1717}{Сумасшедший}{suma'shedshij}{Crazy, mad; a madman}{Ты что, сумасшедший? Как ты можешь ехать на такой скорости?}{Are you crazy? How can you drive at such a speed?}
\vocentry{1718}{Наверняка}{navern’ə'ka}{For sure, certainly}{Я думаю, что сдал тест, но узнаю наверняка только после того, как опубликуют результаты.}{I think I passed the test but I’ll know for sure only when they publish the results.}
\vocentry{1719}{Здорово}{'zdoravə}{Great, awesome}{Должно быть, здорово путешествовать так много, как твой одногруппник.}{That must be great to travel as much as your groupmate does.}
\vocentry{1720}{Напряжение}{napr’ə'zhenije}{Tension; voltage}{Между ними всегда присутствовало некоторое напряжение из-за того, что они раньше встречались.}{There was always some tension between them, because they used to date.}
\vocentry{1721}{Могила}{ma'gila}{A grave}{Эта могила всегда покрыта свежими цветами. Кто здесь похоронен?}{This grave is always covered with fresh flowers. Who’s been buried here?}
\vocentry{1722}{Тоска}{tas'ka}{Melancholy, boredom, sadness}{Твоя тоска убивает тебя! Приободрись и начни действовать!}{Your melancholy is killing you! Cheer up and start acting!}
\vocentry{1723}{Приводить}{priva'dit’}{To bring (with); to lead (to), to result in}{Нельзя приводить с собой на праздник людей, не спросив хозяина.}{One can’t bring people to a party with them without asking the host.}
\vocentry{1724}{Удивляться}{udiv'l’atsa}{To be surprised, to wonder (at)}{Не стоит удивляться, что она такая эгоистичная. Она единственный ребёнок в семье.}{One shouldn’t be surprised that she’s so selfish. She’s an only child in the family.}
\vocentry{1725}{Буквально}{buk'val’nə}{Literally}{Это займёт буквально две минуты. Вам не придётся ждать.}{It’ll take literally two minutes. You won’t have to wait.}
\vocentry{1726}{Борт}{bort}{A board (about a ship or a plane); a side}{Капитан разрешил нам взойти на борт и прогуляться по кораблю.}{The captain allowed us to get on board and to have a walk about the ship.}
\vocentry{1727}{Забрать}{zab'rat’}{To take away}{Есть вещи, которые никто не сможет у нас забрать. Например, любовь и вера.}{There are things that nobody can take away from us. Love and faith, for example.}
\vocentry{1728}{Грех}{greh}{A sin}{Самоубийство – единственный грех, который Бог не прощает.}{Suicide is the only sin that God doesn’t forgive.}
\vocentry{1729}{Несчастный}{ne'shchasnyj}{Unhappy, miserable; unlucky}{Кто сказал, что он несчастный человек? Он просто не умеет ценить то, что имеет.}{Who said he’s an unhappy person? He just can’t appreciate what he has.}
\vocentry{1730}{Строй}{stroj}{A system, an order, a regime}{Политический строй этой страны основан на тирании и насилии.}{The political system of this country is based on tyranny and violence.}
\vocentry{1731}{Занять}{za'n’at’}{To take (about time)}{Полная реконструкция здания должна занять около трёх месяцев.}{The complete reconstruction of the building will take about three months.}
\vocentry{1732}{Исследование}{iss'ledavanije}{Research}{Научное исследование профессора длилось более десяти лет и принесло поразительные результаты.}{The professor’s scientific research lasted more than ten years and brought amazing results.}
\vocentry{1733}{Доказать}{daka'zat’}{To prove}{Мы сможем доказать, что ты невиновен, только если ты будешь говорить правду.}{We’ll be able to prove that you’re innocent only if you tell the truth.}
\vocentry{1734}{Буква}{'bukva}{A letter (alphabet)}{Какая первая буква русского алфавита?}{What’s the first letter of the Russian alphabet?}
\vocentry{1735}{Реакция}{re'aktsija}{A reaction, a response}{Реакция общества на теракт была однозначной.}{The reaction of the society to the terrorist act was unambiguous.}
\vocentry{1736}{Варить}{va'rit’}{To boil, to cook}{Эти креветки не нужно варить. Просто положи их в горячую воду.}{There’s no need to boil the shrimps. Just put them into hot water.}
\vocentry{1737}{Мужской}{muzhs'koj}{Male, masculine; man’s, gentleman’s}{Я услышала незнакомый мужской голос внизу и спустилась проверить, кто это.}{I heard an unfamiliar male voice downstairs and went down to check who it was.}
\vocentry{1738}{Резкий}{'reskij}{Sharp, abrupt}{Резкий спад экономического роста сильно повлиял на зарплаты.}{The sharp decline of the economic growth affected the salaries badly.}
\vocentry{1739}{Пила}{pi'la}{A saw}{Пила застряла в стволе дерева, и я не знаю, как вытащить её.}{The saw is stuck in the tree trunk and I don’t know how to get it out.}
\vocentry{1740}{Литературный}{litera'turnyj}{Literary}{Твой литературный стиль неуместен в наших повседневных разговорах.}{Your literary style is out of place in our everyday talks.}
\vocentry{1741}{Изменить}{izme'nit’}{To change (perfective)}{Мы ничего не можем изменить в этой ситуации. Мы можем только ждать.}{We can’t change anything in this situation. We can only wait.}
\vocentry{1742}{Обезьяна}{abez'jana}{A monkey}{В цирке обезьяна украла мой кошелёк, пока я пытался сфотографироваться с ней.}{In the circus a monkey stole my purse while I was trying to take a photo with it.}
\vocentry{1743}{Бригада}{bri'gada}{A team, a crew (a group of people employed to work together on a task)}{Строительная бригада должна закончить дом к середине ноября, если погода позволит.}{The construction team will have to finish the house by the middle of November if the weather permits.}
\vocentry{1744}{Ловить}{la'vit’}{To catch}{Лягушки могут ловить насекомых своими длинными языками.}{Frogs can catch insects with their long tongues.}
\vocentry{1745}{Остаток}{as'tatək}{A remainder, the rest; a surplus}{Я хочу провести остаток отпуска на даче без суеты и шума.}{I want to spend the remainder of my holiday in the country house without any fuss and noise.}
\vocentry{1746}{Стремиться}{stre'mitsa}{To strive, to aspire}{Для человека естественно стремиться стать лучше.}{It’s natural for a human being to strive for becoming better.}
\vocentry{1747}{Пункт}{punkt}{A point; a station}{Этот пункт твоей презентации лучше заменить чем-то другим.}{It’s better to replace this point of your presentation with something else.}
\vocentry{1748}{Тянуться}{t’ə'nutsa}{To drag on}{Я никогда не думал, что время может тянуться так долго.}{I’ve never thought time could drag on for so long.}
\vocentry{1749}{Локоть}{'lokat’}{An elbow}{Удивительно, но ни один человек не может укусить свой локоть.}{It’s amazing but no man can bite his own elbow.}
\vocentry{1750}{Обязанный}{a'b’azannyj}{Obliged}{Продавец, обязанный принять бракованный товар, отказался это сделать.}{The seller, obliged to accept the faulty good, refused to do it.}
\end{multicols}
